%% LaTeX preamble for Korean Haskell book
%% Uses scrbook (KOMA-Script) + kotex for Korean typesetting

% ── Korean language support ─────────────────────────────────────────────────
\usepackage[hangul]{kotex}          % Korean typography (requires XeLaTeX)

% ── Font configuration ───────────────────────────────────────────────────────
% Uncomment and adjust if specific fonts are installed on your system.
% For broad compatibility the defaults provided by kotex are used.
%
\setmainfont[Ligatures=TeX]{Noto Serif CJK KR}
\setsansfont{Noto Sans CJK KR}
\setmonofont{Hack} % D2Coding, etc

% ── Mathematics ──────────────────────────────────────────────────────────────
\usepackage{amsmath}
\usepackage{amssymb}
\usepackage{amsthm}

\newtheorem{theorem}{정리}[chapter]
\newtheorem{lemma}[theorem]{보조 정리}
\newtheorem{definition}[theorem]{정의}
\newtheorem{example}[theorem]{예제}

% ── Source-code listings ─────────────────────────────────────────────────────
% Quarto renders code blocks via its own highlighting pipeline; these settings
% tune the appearance of the resulting verbatim/fancyvrb output.
\usepackage{setspace}
\usepackage{etoolbox}
\usepackage{fancyvrb}
\usepackage{fvextra}        % line-wrapping inside verbatim
\fvset{baselinestretch=1.0} % 1.0은 기본, 0.8 등으로 줄여서 촘촘하게 가능
\usepackage{mdframed}

% 기존의 잘못된 표지 생성 로직을 완전히 무력화하고 강제로 표지 작성
\makeatletter
\renewcommand{\maketitle}{
  \begin{titlepage}
    \centering
    \vspace*{5cm}
    {\Huge \bfseries \sffamily Haskell 프로그래밍 \par} % 책 제목 직접 명시
    \vspace{2em}
    {\Large \@author \par}
    \vspace{1.5em}
    {\large \@date \par}
  \end{titlepage}
  \thispagestyle{empty}
  \clearpage
  \pagenumbering{roman} % 여기서부터 로마자 번호 시작
}
\makeatother

%% \makeatletter
%% \AtBeginDocument{
%%   % 1. 표지 디자인 실행
%%   \begin{titlepage}
%%     \centering
%%     \vspace*{3cm}
%%     % Quarto가 yml에서 가져온 제목(\@title)을 사용
%%     {\sffamily\bfseries\Huge \@title \par} 
%%     \vspace{1.5cm}
%%     {\Large \@author \par}
%%     \vfill
%%     {\large \@date \par}
%%   \end{titlepage}
%%   
%%   % 2. 본문 습격 방지: 여기서 페이지를 강제로 끊고 초기화
%%   \clearpage
%%   \pagenumbering{arabic}
%%   \setcounter{page}{1}
%%   
%%   % 3. Quarto가 자동으로 호출할 \maketitle을 무력화 (중복 방지)
%%   \let\maketitle\relax 
%% }
%% \makeatother


% Wrap long lines in code blocks instead of overflowing the margin
\DefineVerbatimEnvironment{Highlighting}{Verbatim}{
  breaklines,
  breaksymbolleft={},
  commandchars=\\\{\}
}

% 결과물(verbatim) 환경만 별도로 줄간격과 위아래 여백 조정??
\BeforeBeginEnvironment{verbatim}{\setstretch{1.1}\vspace{-1.7em}} 
\AfterEndEnvironment{verbatim}{\vspace{-.9em}}

% 테두리 스타일 정의
\surroundwithmdframed[
    linecolor=gray!30,       % 테두리 색상
    linewidth=0.5pt,         % 테두리 두께
    % roundcorner=2pt,         % 모서리 둥글게
    skipabove=8pt,         % 위쪽 여백
    % skipbelow=10pt,         % 아래쪽 여백
    innerleftmargin=8pt,    % 왼쪽 여백
    innerrightmargin=8pt,   % 오른쪽 여백
    innertopmargin=4pt,      % 위쪽 여백
    innerbottommargin=4pt,   % 아래쪽 여백
    backgroundcolor=gray!5   % 배경색 (연한 회색)
]{Highlighting}

\newenvironment{stderrbox}
  {\par\medskip\itshape\color{red} \textbf{Error:}\begin{quote}}
  {\end{quote}\medskip}

% ── KOMA-Script chapter/section style ────────────────────────────────────────
\KOMAoptions{
  chapterprefix=true,       % "제 N 장" style headings
  headings=big
}

% ── Hyperlinks (loaded by Quarto; just configure colours here) ───────────────
% \hypersetup must be deferred until after hyperref is loaded by Quarto.
\AtBeginDocument{%
  \hypersetup{
    linkcolor  = black,
    citecolor  = teal,
    urlcolor   = teal
  }%
}

% ── Page geometry ────────────────────────────────────────────────────────────
\usepackage{geometry}
\geometry{
  a4paper,
  top    = 30mm,
  bottom = 30mm,
  inner  = 25mm,
  outer  = 25mm
}

% ── Misc utilities ───────────────────────────────────────────────────────────
\usepackage{microtype}      % microtypographic improvements
\usepackage{booktabs}       % professional-quality tables
\usepackage{graphicx}       % figure inclusion
\usepackage{caption}        % caption customisation
